\documentclass[10pt,a4paper]{amsart}
\usepackage[margin=19mm]{geometry}
\usepackage{amsmath,amssymb,mathtools}
\usepackage[scr=rsfs,cal=euler]{mathalfa}
\usepackage{xfrac}
\usepackage[unicode,hidelinks,bookmarks=false]{hyperref}
\newcommand{\ie}{i.e.\ }
\newcommand{\cf}{cf.\ }
\newcommand{\resp}{respectively\ }
\newcommand{\Subsection}{\subsection}
\providecommand{\nobreakdash}{\nobreak\hbox{-}}
\setlength{\emergencystretch}{2em}
\multlinegap0pt
\usepackage{amssymb}
\usepackage[shortlabels]{enumitem}
\usepackage{xcolor}
\usepackage{dsfont}
\newtheorem{theorem}{Theorem}
\newcommand{\RL}{\mathbb{R}}
\title{Entropic analogues of Gr{\"u}nbaum's inequality}
\author{Matthieu Fradelizi, Lampros Gavalakis, Martin Rapaport}
\date{Source: arXiv:2607.23269v1 (CC BY 4.0)}
\begin{document}
\maketitle

\noindent\emph{Source and licence.} Entropic analogues of Gr{\"u}nbaum's inequality. Version arXiv:2607.23269v1 (CC BY 4.0), distributed by arXiv under the Creative Commons Attribution 4.0 International licence, http://creativecommons.org/licenses/by/4.0/, as recorded in the arXiv licence metadata for this exact version and shown on the abstract page https://arxiv.org/abs/2607.23269v1. Full licence notice: \texttt{licenses/arxiv-2607.23269-CC-BY-4.0.txt}.\par\smallskip

\noindent\textit{Editorial scope.} Counted unit: the article's theorem thm:renyi-truncation (source0.tex lines 784--794). The article's complete original proof of it is delivered with it (source0.tex lines 796--906, one \texttt{proof} environment of 111 lines). No mathematics is written, repaired or re-derived: the statement and the proof are the article's own lines, in the article's own notation, and no retained line is substituted or re-flowed.  Retained uncounted as the source-local context the proof needs: lines 594--603.  Nothing was omitted from any retained span, so the package carries no omission record.  No retained line contains a citation, so the wrapper carries no bibliography and no bibliography entry.  No retained line contains a figure environment, an image-inclusion command or drawing code, so no image is delivered and the package declares no asset dependency.  Editorial formatting only, in addition: an AMS \texttt{amsart} preamble that restates the article's own declarations and macros used by the retained lines, namely the environments \texttt{theorem} on separate counters, together with the standard packages those lines need.  The retained environments print in this document as Theorem 1 in order of appearance, whereas the article prints the same items with the numbering of its own full document.  The complete native source of the article as distributed in the arXiv e-print for this exact version is bundled unchanged as \texttt{sources/206014/source0.tex}.
\begin{theorem}\label{entropygrunbaumTh}
Let $f$ be a log-concave probability density and for $m \in \RL$, let 
\[
f_{m}^-(x):=\frac{f(x)\mathbf 1_{(-\infty,m]}(x)}{\int_{-\infty}^{m} f(y)\,dy}.
\]
Then
\[
h(f_{m}^-)\le h(f).
\]
\end{theorem}

\begin{theorem}%[R\'enyi entropy decreases under truncation]
\label{thm:renyi-truncation}
Let $f$ be a log-concave probability density and for $m \in \RL$, let 
\[
f_{m}^-(x):=\frac{f(x)\mathbf 1_{(-\infty,m]}(x)}{\int_{-\infty}^{m} f(y)\,dy}.
\]
Then, for every $\alpha\in[0,\infty]$, 
\[
h_\alpha\!\big(f_{m}^{-}\big)\le h_\alpha(f).
\]
\end{theorem}

\begin{proof}
The cases $\alpha=0$ and $\alpha=\infty$ were established above.
We may therefore assume that $\alpha\in(0,\infty)$ and, without loss of generality, take $m=0$.
We denote $\mathbb R^- := (-\infty,0]$.  
Define
\[
f^-(x) := \frac{f(x)\mathbf 1_{\mathbb R^-}(x)}{p^{-}},
\qquad
p^{-} := \int_{\mathbb R^-} f(x)\,dx,
\]
and
for $\alpha \ge 0$, let
\[
G(\alpha)
:= \frac{\|f\|_\alpha}{\|f\,\mathbf 1_{\mathbb R^-}\|_\alpha}
= \left(
\frac{\int_{\mathbb R} f(x)^\alpha\,dx}{\int_{\mathbb R^-} f(x)^\alpha\,dx}
\right)^{1/\alpha}.
\]

\medskip

We claim that the function $\alpha \mapsto G(\alpha)$ is non-increasing on $(0,\infty)$.
It suffices to show that $\frac{d}{d\alpha}\log G(\alpha)\le 0$.  
We first write
\[
\log G(\alpha)
=
\frac{1}{\alpha}\log\!\int_{\mathbb R} f(x)^\alpha\,dx
-
\frac{1}{\alpha}\log\!\int_{\mathbb R^-} f(x)^\alpha\,dx.
\]
A direct differentiation yields
\begin{equation}\label{derivada}
\frac{d}{d\alpha}\log G(\alpha)
=
\frac{1}{\alpha^2}
\Bigg[
\alpha\,\frac{\int_{\mathbb R} f(x)^\alpha \log f(x)\,dx}{\int_{\mathbb R} f(x)^\alpha\,dx}
-
\log\!\int_{\mathbb R} f(x)^\alpha\,dx
-
\alpha\,\frac{\int_{\mathbb R^-} f(x)^\alpha \log f(x)\,dx}{\int_{\mathbb R^-} f(x)^\alpha\,dx}
+
\log\!\int_{\mathbb R^-} f(x)^\alpha\,dx
\Bigg].
\end{equation}
To interpret \eqref{derivada} in terms of differential entropies, define the probability densities
\begin{equation} \label{powered}
g_\alpha(x)
:=
\frac{f(x)^\alpha}{\int_{\mathbb R} f(y)^\alpha\,dy},
\qquad
g_\alpha^-(x)
:=
\frac{f(x)^\alpha \mathbf 1_{\mathbb R^-}(x)}{\int_{\mathbb R^-} f(y)^\alpha\,dy}.
\end{equation}
A straightforward computation shows that
\[
h(g_\alpha)
=
-\int_{\mathbb R} g_\alpha(x)\log g_\alpha(x)\,dx
=
-\alpha\,\frac{\int_{\mathbb R} f(x)^\alpha \log f(x)\,dx}{\int_{\mathbb R} f(x)^\alpha\,dx}
+
\log\!\int_{\mathbb R} f(x)^\alpha\,dx,
\]
and similarly,
\[
h(g_\alpha^-)
=
-\int_{\mathbb R^-} g_\alpha^-(x)\log g_\alpha^-(x)\,dx
=
-\alpha\,\frac{\int_{\mathbb R^-} f(x)^\alpha \log f(x)\,dx}{\int_{\mathbb R^-} f(x)^\alpha\,dx}
+
\log\!\int_{\mathbb R^-} f(x)^\alpha\,dx.
\]
Substituting into \eqref{derivada} yields the identity
\begin{equation} \label{Gder}
\frac{d}{d\alpha}\log G(\alpha)
=
-\frac{1}{\alpha^2}
\big(
h(g_\alpha)-h(g_\alpha^-)
\big).
\end{equation}

Since $f$ is log-concave, $f^\alpha$ is log-concave for all $\alpha> 0$, and hence $g_\alpha$ is a log-concave probability density on $\mathbb R$.  
Applying Theorem \ref{entropygrunbaumTh} to $g_\alpha$ and its restriction $g_\alpha^-$, we obtain
\[
h(g_\alpha^-)\le h(g_\alpha).
\]
Together with \eqref{Gder}, this implies
\[
\frac{d}{d\alpha}\log G(\alpha)\le 0,
\]
so that $G(\alpha)$ is non-increasing on $(0,\infty)$.
Consequently,
\[
G(\alpha)\le G(1)\qquad\text{for all }\alpha\ge 1,
\]
and 
\[
G(\alpha)\geq G(1)\qquad\text{for all }\alpha\leq 1,
\]
which is in both cases equivalent to
\[
h_\alpha(f^-)\le h_\alpha(f).
\]
%The case $\alpha = 0$ is straightforward. 
\end{proof}

\end{document}
